\section{Conclusion}\label{sec:conclusion}

The Rule of 40 is a line, not a frontier. Read through a model of the attainable set, it is a price at which growth is traded for margin, and the maximal score is reached where a line at that price touches the frontier. Each firm's frontier is concave, with an intercept set by the margin it earns at market growth and a slope set by the diminishing returns to buying growth. In value terms the Rule of $X$ is a rule of economic profit plus growth at its net value, which carries the frontier beyond software under the conditions of Section~\ref{sec:model-general}. The frontier that can be estimated across firms is drawn through the points firms chose. In SEC-registered software firms from 2010 to 2025, among the best performers, a point of growth has come with well under a point less margin across firms, about a fifth of a point in the baseline and \ScPsiOpTypLo{} to \ScPsiOpTypHi{} of a point for 90\% of the defensible analysis choices, far from the one-for-one trade of the Rule-of-40 line in observed scores, on the organic frontier, under both imputation bounds and in every specification; on the GAAP basis the trade steepened after 2020, to \PsiLateOp, with an interval that excludes zero in \PsiDiffOpCountNeg{} views, but that interval depends on the period split and the sample. Between 2016--19 and 2023--25 the 95th percentile of the score fell by \TopFallOpEone{} points on the GAAP basis and \TopFallCashEone{} after stock-based compensation in the baseline, and by \ScTopOpCoreFallLo{} to \ScTopOpCoreFallHi{} and \ScTopCashCoreFallLo{} to \ScTopCashCoreFallHi{} points across the sample choices, while fewer firms grew fast. The GAAP decline survives a familywise correction at 10\% under \TopFamTen{} of four procedures, and at 5\% under \TopFamFive, and the cash decline under none, the data do not separate a change in the growth mix from a fall in the best margins of fast growers, and the top partly recovered in 2025 under observed scores. The data do not identify the weight on growth that firms act on, the curvature of the growth technology or a decline of the score with scale, and the lower cash scores of capital-intensive firms are an accounting effect.

The answer to ``Rule of 100 or Rule of 1000?'' is neither, as a steady state. Automation raises the margin intercept by removing operating cost, but it adds computing cost and the recurring costs of maintaining, securing and supervising AI, and, for firms that own capacity, capital that weighs on cash-basis scores. Saturation removes the growth surplus as markets fill, and within one market a finite budget of share headroom limits how long any surplus lasts. Under free entry into a market for a homogeneous product, competition removes the margin that automation creates unless the firm holds a scarce input, and even a rent cannot lift the score above the full-cost ceiling. A sustained cash-basis score above 100 therefore requires, at the full-automation cost structure of Section~\ref{sec:horizons-ceiling} and with negligible capital needs, a market growing at a rate few markets sustain: about \GstarCenOne\% a year if compute absorbs \CcCen\% of revenue at 2025 prices and the other costs are at their central values, between \GstarCcLo{} and \GstarCcHi\% over the scenario range of that share, and \GstarPhiHalfOne{} or \GstarPhiSevenOne\% if half or 70\% of cost savings pass to customers. The threshold falls for a firm with negative cash capital or costs leaner than that structure, and rises with positive capital. Reaching that cost structure is the scenario's route, near-complete automation; \autoref{cor:ceiling}(c) notes that a partial program can do better when some processes do not pay. Scores of several hundred are possible only briefly, in new markets with exceptional demand pull. Capacity to execute risky transitions limits how fast a firm gets to its ceiling: under our scenario assumptions a top-tail firm realizes about \PaceCenSix\% of the benefits of its automation program by 2036 under central execution, having embedded about \CenEmbSix\% of it, and its upper bound reaches all of it by \PaceFastDone{} under fast execution. On the central scenario path, at fixed prices, its maximal durable score is about \CenREight{} in 2028, \CenROne{} in 2031 and \CenRSix{} in 2036 (\CenRSixNet{} net of the dip of transitions in flight; \CenRSixLo{} to \CenRSixHi{} from anchor sampling alone), and in 2036 \CenRSixAllSM{} if all and \CenRSixNoSM{} if none of its selling expense buys growth. The program repays the costs of its transitions only after \PayBackCen{} years, so under central execution it does not pay back within five years except at a few points of the joint range of the other inputs. These are scenarios, not forecasts.

For practitioners the implications are concrete. Compute the score on annual free cash flow after stock-based compensation and capital expenditure, state the basis, since scores on this basis are not comparable with the ones practitioners usually quote, and charge acquired growth with what it cost, charging either interest or a cost of capital on the capital that paid for it, not both. This paper does not test the Rule of 40 as an investment screen. As a rule for setting growth spending it is the wrong test: a firm can accept a lower score while its marginal growth spend earns more than it costs, which it must show from the payback and retention of its customer cohorts at a stated discount rate; the data do not show which firms meet that test. Read cross-sectional comparisons as a record of choices by firms with different opportunities, not as the trade-off any one firm faces. Supplement the score with economic profit (\autoref{prop:general}) and gross profit per dollar of computing cost as headcount stops being the main cost. And evaluate AI strategies by the parameters they move: the margin intercept, demand pull, capital intensity, pricing power, the accountability floor, the full-cost floors of maintaining, securing and supervising AI, and the capacity to execute the transitions that get the firm there. The gains from a better growth engine are transient and those from cutting overhead and from retention durable, so the growth engine matters most where a firm has share to gain and overhead where it has little, while retention counts at every stage (\autoref{cor:priorities}), and which comes first at a given position depends on the weight the owner acts on and on how much selling spend buys growth: net of maintenance selling spend, cutting general and administrative cost outranks automating the growth engine at \lpNGACutFirstNet{} of \lpNPositions{} observed positions at the central maintenance level, and the companion paper ranks engineering above both (Section~\ref{sec:levers-position}); savings in cost of revenue are worth only what pricing power lets the firm keep; owned capacity pays only where growth is slow enough for the cost it removes to exceed growth times the capital it adds; and control designs that rely on one model leave a requirement that more layers do not meet. Under free entry into a homogeneous market, automation's savings raise the long-run score through a rent only for a firm that holds an input that stays scarce, and then only up to the full-cost ceiling; a firm without one keeps a normal return on its capital, not a rent.
